% Please keep the comparison's goal, obstacle, and bridge on the same page.
\begin{samepage}
\subsection{Adaptive midpoint comparison and instance prior}
\label{app:lower-proof-midpoint}

We relate the target law to the sampler's output law through the discrepancies
at individual commits. Three points guide the comparison:
\begin{itemize}
\item \emph{Goal.} Express the output affinity, or equivalently one minus the
squared-Hellinger error, in terms of local commit affinities.
\item \emph{Obstacle.} A different committed value can change subsequent
probes, batches, and the stopping round. Thus local affinities depend on the
realized history; the fixed-product affinity formula does not directly
combine them.
\item \emph{Bridge.} Normalize the geometric mean of the two commit kernels
at each history. These normalized kernels define a common path law under
which the product of local affinities equals the output affinity in
expectation (\cref{app:lem:adaptive-midpoint-lower}).
\end{itemize}
\end{samepage}

\emph{Fixed objects.}
Fix an admissible algorithm with finite pathwise budgets
\((\QueryBudget{},\RoundBudget{})\), an instance
\(I=(\HiddenMatching,\TriggerVector)\) of
\cref{app:def:hard-matching-family}, and one value \(w\) of the
decision-rule seed. Put
\begin{equation}
 \HardOutputLaw{I}{w}
 :=\SamplerOutputLaw{w}{\HardOracle{I}}.
 \label{app:eq:hard-seeded-output-law}
\end{equation}
The following three executions differ in their commit kernels, defined
below. The exact comparison process has terminal law \(\HardTargetLaw{I}\);
Step~2 of the proof verifies this even for adaptive batches.

\begin{center}
\small
\begin{tabularx}{\linewidth}{@{}>{\raggedright\arraybackslash}p{0.23\linewidth}
  >{\raggedright\arraybackslash}p{0.32\linewidth}
  >{\raggedright\arraybackslash}X@{}}
\toprule
Execution & Commit kernel at history \(h\) & Law and role \\
\midrule
Exact comparison
& \(\HardExactBatchLaw{t}{I}(\cdot\mid h)\): true joint conditional
& \(\HardTargetLaw{I}\): target output law. \\
\addlinespace
Product-commit sampler
& \(\HardProductBatchLaw{t}{I}(\cdot\mid h)\): product of oracle rows
& \(\HardOutputLaw{I}{w}\): output law being evaluated. \\
\addlinespace
Adaptive midpoint
& \(\HardMidpointKernel{t}{I}(\cdot\mid h)\): normalized geometric mean
& \(\HardMidpointLaw{I}{w}\): analytic path law for averaging local affinities. \\
\bottomrule
\end{tabularx}
\end{center}

\begin{itemize}
\item \emph{Shared operations.} All three use the same seed-fixed decision
rule and the same frozen oracle \(\HardOracle{I}\) for probes. At a common
history they select the same batch \(B_t(h)\).
\item \emph{Changed operation.} Only the law of the committed values changes.
The resulting histories, and hence later batches, can differ between executions.
\end{itemize}

\begin{definition*}[Commit kernels and adaptive midpoint law]
\emph{History and batch.}
Immediately before commit \(t\), after the intervening probes, let \(h\)
be the complete observable history: committed positions and values,
the earlier oracle transcript, and the decision-rule state.
Thus $h=\OpVisibleHistory{s_t}$ in the operation notation of
\cref{app:eq:operation-action,app:eq:operation-visible-history}.
Write \(G(h)\) for the committed set, \(x_{G(h)}\) for its values, and
\(y_t(h)\) for the commit state in \eqref{app:eq:commit-state}.
For fixed \((I,w)\), the probes and batch selection are deterministic
functions of the observed past; write \(B_t(h)\) for the selected batch.
The seed \(w\) is suppressed from local kernels because \(h\) includes
the seed-fixed decision-rule state, but retained in full-law and
terminal-round notation.

\emph{Commit kernels.}
For \(z\in\Vocab^{B_t(h)}\), define
\begin{align}
 \HardExactBatchLaw{t}{I}(z\mid h)
 &:=\HardTargetLaw{I}\bigl(
   X_{B_t(h)}=z\mid X_{G(h)}=x_{G(h)}
  \bigr),
 \label{app:eq:hard-exact-batch-law}\\
 \HardProductBatchLaw{t}{I}(z\mid h)
 &:=\prod_{j\in B_t(h)}\HardOracleRowMass{I}{y_t(h)}{j}{z_j}.
 \label{app:eq:hard-product-batch-law}
\end{align}

\emph{Local affinity and midpoint law.}
Define
\begin{align}
 \HardLocalAffinity{t}{I,h}
 &:=\sum_{z\in\Vocab^{B_t(h)}}
    \sqrt{
     \HardExactBatchLaw{t}{I}(z\mid h)
     \HardProductBatchLaw{t}{I}(z\mid h)},
 \label{app:eq:hard-local-affinity}\\
 \HardMidpointKernel{t}{I}(z\mid h)
 &:=\frac{
    \sqrt{
     \HardExactBatchLaw{t}{I}(z\mid h)
     \HardProductBatchLaw{t}{I}(z\mid h)}}
   {\HardLocalAffinity{t}{I,h}}.
 \label{app:eq:hard-midpoint-kernel}
\end{align}
Let \(\HardMidpointLaw{I}{w}\) be the analytic adaptive path law that runs
the same decision rule and frozen-oracle probes, drawing every commit
batch from \(\HardMidpointKernel{t}{I}(\cdot\mid h)\).
\end{definition*}

\paragraph{Why the construction is well defined.}
\begin{itemize}
\item \emph{Conditioning on the history.}
For fixed \((I,w)\), every reachable \(h\) can be reconstructed from
\(x_{G(h)}\): replay the seed-fixed decision rule from the initial state and,
at each earlier commit, supply the coordinates of \(x_{G(h)}\) selected
by that commit. The frozen oracle reproduces the intervening replies,
batches, and decision-rule states. Conversely, \(h\) records \(x_{G(h)}\).
Thus conditioning on this \(h\) adds no random event beyond
\(X_{G(h)}=x_{G(h)}\), justifying
\eqref{app:eq:hard-exact-batch-law}.
\item \emph{Normalization.}
Strict positivity makes the denominator nonzero. Cauchy--Schwarz gives
\(0<\HardLocalAffinity{t}{I,h}\le1\). In particular,
\[
 \sum_{z\in\Vocab^{B_t(h)}}\HardMidpointKernel{t}{I}(z\mid h)
 =\frac{\displaystyle
   \sum_{z\in\Vocab^{B_t(h)}}
    \sqrt{\HardExactBatchLaw{t}{I}(z\mid h)
           \HardProductBatchLaw{t}{I}(z\mid h)}}
  {\HardLocalAffinity{t}{I,h}}
 =\frac{\HardLocalAffinity{t}{I,h}}{\HardLocalAffinity{t}{I,h}}
 =1.
\]
\item \emph{Pathwise budgets.}
Both commit kernels are strictly positive, so the midpoint kernel has
the same support at every reachable history. It therefore traverses the
same decision tree as the implemented sampler. Every midpoint path has
at most \(\QueryBudget{}\) counterfactual submissions and at most
\(\RoundBudget{}\) nonempty commits.
\end{itemize}

\paragraph{From local affinity to output error.}
The defining factorization is
\[
 \sqrt{\HardExactBatchLaw{t}{I}(z\mid h)
        \HardProductBatchLaw{t}{I}(z\mid h)}
 =\HardLocalAffinity{t}{I,h}\HardMidpointKernel{t}{I}(z\mid h).
\]
Along a complete path, the normalized \(\HardMidpointKernel{t}{I}\) factors
supply its midpoint probability, while the \(\HardLocalAffinity{t}{I,h}\)
factors retain the local affinity losses.
The next identity averages their product over the adaptive paths.

\begin{lemma}[Adaptive midpoint identity]
\label{app:lem:adaptive-midpoint-lower}
For every fixed instance \(I\) and decision-rule seed \(w\),
\begin{equation}
 1-\sqHellinger(
      \HardTargetLaw{I},\HardOutputLaw{I}{w})
 =\mathbb E_{H\sim\HardMidpointLaw{I}{w}}
   \prod_{t=1}^{\HardTerminalRoundCount{I}{w}{H}}
    \HardLocalAffinity{t}{I,H_{t-1}},
 \label{app:eq:adaptive-midpoint-identity}
\end{equation}
where \(\HardTerminalRoundCount{I}{w}{H}\le \RoundBudget{}\) is the terminal number of
nonempty commit batches on the midpoint path.
Here \(H_{t-1}\) denotes the complete observable history immediately
before commit \(t\), after all intervening counterfactual operations.
The subscript counts earlier commits, not individual query operations.
\end{lemma}

The remaining lower-bound argument uses this identity at two levels:
\begin{itemize}
\item \emph{Local penalty.} \Cref{app:lem:collision-affinity} bounds the
factors \(\HardLocalAffinity{t}{I,h}\) using unresolved edges whose
endpoints are committed together.
\item \emph{Global error.} The survivor and collision arguments count
these events under the instance--midpoint law
\eqref{app:eq:joint-instance-midpoint-law}. The identity converts that
count into an output-error bound.
\end{itemize}

\begin{proof}

\emph{1. Reconstruct the unique adaptive path.}
Fix a complete output \(x\in\Vocab^N\).  Starting with the empty history,
replay the deterministic seed-fixed decision rule.  Whenever it chooses
\(B_t(h_{t-1}(x))\), feed it
\(z_t(x):=x_{B_t(h_{t-1}(x))}\).  This recursively determines the histories,
batches, and terminal round \(\HardTerminalRoundCount{I}{w}{x}\).  No
remasking makes the realized batches disjoint and exhaustive.  Put
\(G_t(x):=G_{t-1}(x)\cup B_t(h_{t-1}(x))\), with \(G_0(x)=\varnothing\).

\emph{2. Factor the exact and product output masses.}
The exact conditional factors telescope:
\begin{align*}
 &\prod_{t=1}^{\HardTerminalRoundCount{I}{w}{x}}
   \HardExactBatchLaw{t}{I}(z_t(x)\mid h_{t-1}(x))\\
 &\quad=\prod_{t=1}^{\HardTerminalRoundCount{I}{w}{x}}
   \frac{\HardTargetLaw{I}(X_{G_t(x)}=x_{G_t(x)})}
        {\HardTargetLaw{I}(X_{G_{t-1}(x)}=x_{G_{t-1}(x)})}
 =\frac{\HardTargetLaw{I}(X_{\PositionSet}=x)}{1}
 =\HardTargetLaw{I}(x).
\end{align*}
Here the numerator at round \(t\) is the denominator at round \(t+1\);
the initial event has probability one and the terminal set is \(\PositionSet\).
This identity and the
product-commit rule therefore give
\begin{align}
 \HardTargetLaw{I}(x)
 &=\prod_{t=1}^{\HardTerminalRoundCount{I}{w}{x}}
   \HardExactBatchLaw{t}{I}
    (z_t(x)\mid h_{t-1}(x)),
 \label{app:eq:hard-exact-path-factorization}\\
 \HardOutputLaw{I}{w}(x)
 &=\prod_{t=1}^{\HardTerminalRoundCount{I}{w}{x}}
   \HardProductBatchLaw{t}{I}
    (z_t(x)\mid h_{t-1}(x)).
 \label{app:eq:hard-product-path-factorization}
\end{align}
The batch-selection rule has no likelihood factor because it is deterministic
given the displayed history.


\emph{3. Normalize the midpoint path mass.}
For the same terminal output, the midpoint path mass is
\begin{equation}
 \HardMidpointLaw{I}{w}(x)
 :=\prod_{t=1}^{\HardTerminalRoundCount{I}{w}{x}}
   \HardMidpointKernel{t}{I}
    (z_t(x)\mid h_{t-1}(x)).
 \label{app:eq:hard-midpoint-path-mass}
\end{equation}
To verify that this path mass sums to one, first account for paths that
finish at different rounds.  Nonempty batches are disjoint, so every path
has at most \(N\) rounds.  For this calculation only, pad each finished path
to \(N\) rounds by setting \(B_t(h)=\varnothing\),
\(\Vocab^\varnothing=\{\varnothing\}\), and
\(\HardMidpointKernel{t}{I}(\varnothing\mid h)=1\) after termination.
These extra factors leave \eqref{app:eq:hard-midpoint-path-mass} unchanged.
Each complete output corresponds to exactly one padded path, and conversely,
because the seed-fixed decision rule is deterministic and the genuine batches
are disjoint and exhaustive.

Write \(h_{t-1}\) for the history determined by the preceding assignments
\(z_1,\ldots,z_{t-1}\).  For any \(1\le r\le N\), hold
\(z_1,\ldots,z_{r-1}\) fixed.  The earlier factors do not depend on \(z_r\),
so the local normalization above gives
\begin{align*}
 &\sum_{z_r\in\Vocab^{B_r(h_{r-1})}}
   \prod_{t=1}^{r}\HardMidpointKernel{t}{I}(z_t\mid h_{t-1})\\
 &\quad=
   \left[\prod_{t=1}^{r-1}
     \HardMidpointKernel{t}{I}(z_t\mid h_{t-1})\right]
   \underbrace{\sum_{z_r\in\Vocab^{B_r(h_{r-1})}}
     \HardMidpointKernel{r}{I}(z_r\mid h_{r-1})}_{=1}\\
 &\quad=\prod_{t=1}^{r-1}
     \HardMidpointKernel{t}{I}(z_t\mid h_{t-1}).
\end{align*}
Applying this equality for \(r=N,N-1,\ldots,1\) removes the nested sums
from the innermost one outwards and ends with the empty product \(1\):
\[
 \sum_{x\in\Vocab^N}\HardMidpointLaw{I}{w}(x)
 =\sum_{z_1\in\Vocab^{B_1(h_0)}}\cdots
  \sum_{z_N\in\Vocab^{B_N(h_{N-1})}}
   \prod_{t=1}^{N}\HardMidpointKernel{t}{I}(z_t\mid h_{t-1})
 =1.
\]

\emph{4. Convert the affinity sum into a midpoint expectation.}
Substituting
\(\sqrt{\HardExactBatchLaw{t}{I}\HardProductBatchLaw{t}{I}}
 =\HardLocalAffinity{t}{I,h}\HardMidpointKernel{t}{I}\) into the terminal
affinity sum yields
\begin{align*}
 &\HellingerAffinity(
    \HardTargetLaw{I},\HardOutputLaw{I}{w})\\
 &\quad=\sum_{x\in\Vocab^N}
  \sqrt{\HardTargetLaw{I}(x)\HardOutputLaw{I}{w}(x)}\\
 &\quad=\sum_{x\in\Vocab^N}
  \HardMidpointLaw{I}{w}(x)
  \prod_{t=1}^{\HardTerminalRoundCount{I}{w}{x}}
   \HardLocalAffinity{t}{I,h_{t-1}(x)}\\
 &\quad=\mathbb E_{H\sim\HardMidpointLaw{I}{w}}
  \prod_{t=1}^{\HardTerminalRoundCount{I}{w}{H}}
   \HardLocalAffinity{t}{I,H_{t-1}}.
\end{align*}
Since \(\HellingerAffinity(p,q)=1-\sqHellinger(p,q)\), this is
\eqref{app:eq:adaptive-midpoint-identity}.
\end{proof}

\paragraph{Instance prior for the remaining proof.}
Keep the seed \(w\) fixed. Draw an instance
\(I=(\HiddenMatching,\TriggerVector)\) from the prior
\(\HardInstancePrior\) as follows:
\begin{itemize}
\item \emph{Matching.} \(\HiddenMatching\) is uniform over all perfect
matchings of \(\PositionSet\).
\item \emph{Triggers.} \(\TriggerIndex{1},\ldots,\TriggerIndex{N}\) are
independent uniform elements of \([\HardBankSize]\), independent of
\(\HiddenMatching\).
\end{itemize}
All random variables used in the remaining proof are evaluated under the
joint analytic law
\begin{equation}
 I\sim\HardInstancePrior,
 \qquad
 H\mid(I,w)\sim\HardMidpointLaw{I}{w}.
 \label{app:eq:joint-instance-midpoint-law}
\end{equation}
